% ICASSP 2027 manuscript — FORMAT SKELETON ONLY (2026-09-18).
% Uses the official kit in this directory: spconf.sty + IEEEbib.bst (downloaded 2026-09-18).
% The section layout mirrors the plan v2 (docs/PLAN_spiking_eegnex_icassp2027.md §6) as ONE candidate structure.
% Title, framing, which experiments/tables go in, and all prose are decided by the senior author; nothing here is
% adopted until then. No result number has been filled in.
% RULES (see AGENTS.md §7): every number traces to one CSV row under outputs/analysis/ or
% results_snapshot_2026-09-13/**/tables/; two main tables are never mixed; no energy multiplier in the title;
% forbidden phrasings listed in .cursor/rules/40-paper.mdc.
% Placeholders are marked \todo{...}; remove the \todo macro output before submission.
% --------------------------------------------------------------------------
\documentclass{article}
\usepackage{spconf,amsmath,graphicx,hyperref}
\usepackage{booktabs,multirow}
\usepackage{xcolor}
\newcommand{\todo}[1]{\textcolor{red}{[TODO: #1]}}

% Notation
\def\x{{\mathbf x}}
\def\R{{\mathbf R}}

% Title: all caps is applied by the style file.
% 题目、框架、选用哪些实验 —— 由学长决定后再填。下面只是模板格式骨架，不代表任何已定的论文结论。
\title{\todo{TITLE — to be decided by the senior author}}

% Single-blind: real names and affiliations. Multiple affiliations: use \twoauthors (see official_kit/Template.tex).
\name{\todo{Author list in agreed order}\thanks{\todo{Funding.} Portions of the code and text were drafted with the assistance of an AI coding agent (Cursor); all experiments, analyses and claims were verified by the authors.}}
\address{\todo{Affiliation(s), city, country}}

\begin{document}
%\ninept
\maketitle

\begin{abstract}
% 100–150 words. Write only after the senior author has fixed the story and the experiment set.
% Candidate framing from the plan v2 (docs/PLAN_spiking_eegnex_icassp2027.md §0.3) is NOT adopted here.
\todo{Abstract — after the senior author decides the framing and which experiments go in.}
\end{abstract}

\begin{keywords}
\todo{Keywords — after the framing is decided}
\end{keywords}

% --------------------------------------------------------------------------
\section{Introduction}
\label{sec:intro}
% Argument (PLAN §1): (1) variance is in subjects, not models — own evidence (docs/PRIOR_RESULTS.md §3.1):
%   top-8 within 1.0 pp, SD 12 pp, 67 M-param DINO models at the bottom, 3 k-param EEGNet within 1 pp of the best,
%   train–test gap 15–33 pp, domain-adversarial training is a net loss on strong backbones.
% (2) Why EEGNeX: learnable FBCSP + multi-scale envelope; first or tied-first on all three cohorts, 58 k params.
% (3) SNN promise (event-driven, sample-synchronous) and SNN fragility (threshold is scale-sensitive -> subject shift
%   moves firing rates) -> label-free alignment as homeostasis; the interaction hypothesis.
% Relation to prior work (mandatory): Zhang et al. 2025 (PLIF on raw samples + EA), NiSNN-A (54-subject LOSO +
%   energy), HR-SNN, conversion methods (LENet/SRIF). What is new here: first per-subject-paired, statistically
%   pre-registered comparison of an SNN against 18 strong ANNs on 87 subjects; quantified alignment x network-type
%   interaction; energy accounting from measured firing rates with EEGNet as reference.
\todo{Introduction — after the senior author decides the framing (≈0.8 column). The comments above summarise the
plan-v2 argument only as background material; do not treat it as the adopted story.}

\noindent\textbf{Contributions.} \todo{To be decided by the senior author (the plan v2 proposed three; see docs/PLAN\_spiking\_eegnex\_icassp2027.md §0.3).}

% --------------------------------------------------------------------------
\section{Method}
\label{sec:method}

\subsection{From EEGNeX to S-EEGNeX}
\label{ssec:mapping}
% Layer mapping (PLAN §2.1): conv1 BN LIF1 | conv2 BN | dw BN pool LIF2 drop | conv3 BN | conv4 BN pool LIF3 drop | fc.
% Weight shapes identical to EEGNeX (58,770 params + a few tau). Readout = non-spiking leaky integrator (LI) of the
% flattened spike-count map (binary readout is an ablation). Fig. 1 here.
\begin{figure}[t]
  \centering
  \todo{Fig.~1: layer mapping EEGNeX $\rightarrow$ S-EEGNeX; D (T=4 direct) vs S (temporal-native 1000/250/63 steps).}
  % \includegraphics[width=\linewidth]{figures/fig1_layer_mapping}
  \caption{Layer-wise spiking of EEGNeX. Left: EEGNeX. Right: S-EEGNeX-S, where LIF$_1$ runs at the 250~Hz sample rate and later neuron layers at the pooled rates; LIF layers are placed after the (linear) average pooling so that all downstream layers receive binary input.}
  \label{fig:mapping}
\end{figure}
\todo{Text for the mapping and the two encodings (D: direct, $T{=}4$; S: temporal-native). Say ``sample-synchronous; with `same' padding the look-ahead is $\le 128$~ms''. Do NOT write ``anytime/streaming''.}

\subsection{Temporal-native encoding}
\label{ssec:encoding}
% LIF1 on the raw sample axis rectifies at threshold; conv2 smooths the non-negative spike train -> envelope estimate
% -> LIF2. State as hypothesis (supported by the analog-front-end ablation), not as a theorem.
\todo{LIF dynamics: $H[t]=V[t-1]+(X[t]-V[t-1])/\tau$, $S=\Theta(H-v_{th})$; ATan surrogate; hard reset in the encoder, soft reset in deeper layers.}

\subsection{Heterogeneous learnable time constants (HetPLIF)}
\label{ssec:hetplif}
% Per-channel learnable tau, log-uniform init: encoder [1.5,4] steps, deeper [1.5,8] steps. Relation to EEGNeX's two
% dilated convolutions (fixed time scales) -> a continuum of learnable time scales. H3: LIF <= PLIF <= HetPLIF.
\todo{One paragraph.}

\subsection{Label-free subject alignment}
\label{ssec:alignment}
% EA per session (He & Wu 2020; 10% shrinkage), then per-trial channel z-score (chosen on the SHU pilot by validation
% macro-F1, D-021). AdaBN: BN statistics re-estimated on the held-out subject's unlabeled trials at evaluation time.
% Both are transductive-unsupervised; the ANN references receive the same alignment arm.
\todo{One paragraph; state transductive use explicitly.}

\subsection{Energy accounting}
\label{ssec:energy}
% E_SNN = sum_analog MACs*4.6 pJ + sum_spike MACs*r*0.9 pJ (45 nm, Horowitz 2014; Dampfhoffer 2023), r = measured
% firing rate per test subject. Reference = EEGNet ANN and S-EEGNet.
\todo{Two sentences + the equation.}

% --------------------------------------------------------------------------
\section{Experiments}
\label{sec:exp}

\subsection{Data, protocol and statistics}
\label{ssec:protocol}
% Cohorts (docs/PRIOR_RESULTS.md §1): WBCIC-SHU 2C 51 subj / 58 ch; SHU 2022 25 subj / 32 ch; WBCIC-SHU 3C 11 subj LOSO.
% All 250 Hz, 0–4 s. 5-fold subject k-fold 6:2:2 with split seeds 41/42 (2C, SHU); LOSO (3C). Recipe: Adam 1e-4,
% batch 128, patience 25 on val macro-F1, single training seed. Same folds as the 18-entry ANN benchmark -> pairing.
% Pre-registered family (2C, Holm, 5 tests): H0-NI (delta=1 pp), H1 primary (difference of differences), H2-NI (S vs D),
% H3 (Page's L), B1 (S-EEGNeX vs S-EEGNet). SHU/3C replication; SHU primary metric balanced accuracy.
\todo{Write (≈0.6 column).}

\subsection{Main comparison against published ANN baselines (unaligned)}
\label{ssec:main}
% ------------------------------------------------------------------------------------------------
% 2026-09-18 (senior's request via WeChat): table restricted to the models we actually compare against,
% each with its publication year; format of the table unchanged. Rows are GENERATED, not typed:
%   python scripts/build_comparison_table.py --out outputs/analysis/comparison_table_2026-09-18
%   -> outputs/analysis/comparison_table_2026-09-18/tab_main_rows.tex   (paste between \midrule ... \bottomrule)
% Sources: 2C/SHU = results_snapshot_2026-09-13/experiments/{wbci_shu/wbci_2c,shu/shu}_cross_subject_bench_v1/tables/per_model.csv
%          3C     = results_snapshot_2026-09-13/experiments/wbci_shu/paper_baseline_3c_821_v1/tables/per_cell.csv (ddof=1)
%          SNN    = outputs/analysis/snn_round/shu_pilot_lr1e4/{per_model,energy}.csv (SHU pilot, split 41 only)
% When the main run finishes, re-run the script with --snn-2c/--snn-shu/--snn-3c pointing at the new
% analyze_snn_round.py outputs and replace the rows; the dagger marks then disappear automatically.
% Twelve in-house entries of the 18-entry benchmark (S4ERP*, DualCD*, ATCNet-S4*, ATCNet ablations) are
% deliberately NOT in this table (report only; they may still enter the CD diagram, k = 19).
% Senior (09-18): "到时候不对比这个 eegnex" — if the EEGNeX row is finally dropped, delete ONE line below;
% kept for now because the spiking models are derived from it (see docs/reports/REPORT_comparison_table_2026-09-18.md §5.4).
% Candidate 2025/2026 baselines requested by the senior (not run yet; choose, then add a row + run):
%   Zhang et al. 2025 (Neural Networks 191:107741, lightweight PLIF SNN, code github.com/Zhr1110/SnnForMI) \cite{zhang2025snn_mi_ea}
%   CBraMod (ICLR 2025, EEG foundation model; weights already in ../Benchmark/weights/CBraMod.pth) \cite{wang2025cbramod}
%   MIRepNet (arXiv 2025, MI-specific foundation model, HUST; pretraining corpus does not include SHU/WBCIC-SHU) \cite{liu2025mirepnet}
%   DSGNet (IEEE JBHI 2026; architecture preview only, training code unreleased) \cite{dsgnet2026}
% ------------------------------------------------------------------------------------------------
\begin{table}[t]
  \centering
  \caption{Cross-subject accuracy (\%, mean $\pm$ SD over subjects, best-validation checkpoint) of published ANN baselines and the proposed spiking models on the three cohorts. Unaligned block: all entries use per-trial z-scoring only and the identical recipe and subject folds. Params at 58 channels / 2 classes. Energy: 45~nm estimate (MAC 4.6~pJ, AC 0.9~pJ, measured firing rates), SHU 32-channel 4-s trial; ``--'' = not computed. $^\dagger$SHU pilot, split seed 41 only (each subject tested once); main run (seeds 41+42) in progress. $^\ddagger$3C from the LOSO paper-recipe run (patience 100); being re-run with the common recipe.}
  \label{tab:main}
  \small
  \begin{tabular}{lccccc}
    \toprule
    Model (year) & 2C (51) & SHU (25) & 3C (11) & Params & $\mu$J/trial \\
    \midrule
    DeepConvNet~\cite{schirrmeister2017deep} (2017) & 82.8 $\pm$ 12.1 & 61.0 $\pm$ 10.3 & -- & 302.6k & -- \\
    EEGNet~\cite{lawhern2018eegnet} (2018) & 82.5 $\pm$ 12.0 & 61.5 $\pm$ 10.7 & 66.4 $\pm$ 13.1$^\ddagger$ & 3.0k & 78 \\
    FBCNet~\cite{mane2021fbcnet} (2021) & 70.0 $\pm$ 10.8 & 59.2 $\pm$ 6.7 & -- & 20.7k & -- \\
    ATCNet~\cite{altaheri2023atcnet} (2023) & 83.3 $\pm$ 12.0 & 62.0 $\pm$ 11.2 & 68.4 $\pm$ 13.3$^\ddagger$ & 114.6k & -- \\
    EEGNeX~\cite{chen2024eegnex} (2024) & 83.4 $\pm$ 11.5 & 65.5 $\pm$ 10.0 & 69.9 $\pm$ 13.6$^\ddagger$ & 58.8k & 2539 \\
    EEG-Deformer~\cite{ding2025deformer} (2025) & 82.7 $\pm$ 11.9 & 61.4 $\pm$ 9.5 & 68.5 $\pm$ 14.8$^\ddagger$ & 1.61M & -- \\
    \midrule
    S-EEGNet (PLIF) & \todo{} & 60.6 $\pm$ 11.1$^\dagger$ & \todo{} & 3.0k & 79 \\
    S-EEGNeX-D (HetPLIF, $T{=}4$) & \todo{} & 64.2 $\pm$ 10.1$^\dagger$ & \todo{} & 58.9k & 358 \\
    \textbf{S-EEGNeX-S (HetPLIF, temporal-native)} & \todo{} & 63.3 $\pm$ 11.0$^\dagger$ & \todo{} & 58.9k & 148 \\
    \bottomrule
  \end{tabular}
\end{table}
\begin{figure}[t]
  \centering
  \todo{Fig.: critical-difference diagram, 2C, k=19 (18 ANN + S-EEGNeX-S), Bonferroni–Dunn vs S-EEGNeX.}
  \caption{Critical-difference diagram on WBCIC-SHU 2C (Friedman + Nemenyi; $N{=}51$ subjects).}
  \label{fig:cd}
\end{figure}
\todo{Non-inferiority statement with one-sided bootstrap lower bound of the paired difference vs $\delta = 1$~pp; Hodges–Lehmann shift, BCa CI, win/tie/loss, $n_{\mathrm{eff}}$.}

\subsection{Alignment $\times$ network type}
\label{ssec:interaction}
\begin{table}[t]
  \centering
  \caption{Aligned block (2C): per-subject accuracy of three models under no alignment, EA, and EA+AdaBN ($3\times 3$). The pre-registered H1 is the difference-of-differences test on the $\{$EEGNeX, S-EEGNeX-S$\}\times\{$none, EA+AdaBN$\}$ subset. \todo{Fill from interaction\_H1.csv and per\_model.csv (run analyze\_snn\_round.py with --aligned-arm eaz).}}
  \label{tab:align}
  \small
  \begin{tabular}{lccc}
    \toprule
    Model & none & EA & EA+AdaBN \\
    \midrule
    EEGNeX & \todo{} & \todo{} & \todo{} \\
    EEGNet & \todo{} & \todo{} & \todo{} \\
    S-EEGNeX-S & \todo{} & \todo{} & \todo{} \\
    \bottomrule
  \end{tabular}
\end{table}
\begin{figure}[t]
  \centering
  \todo{Fig.: per-subject alignment gain, ANN vs SNN (paired), with 95\% CI of the difference of differences.}
  \caption{Alignment gain per subject for EEGNeX and S-EEGNeX-S on 2C.}
  \label{fig:interaction}
\end{figure}
\todo{Report H1 as the primary test; if not significant, report as an observation (D-023).}

\subsection{Firing rates, energy and ablations}
\label{ssec:ablation}
\begin{figure}[t]
  \centering
  \todo{Fig.: between-subject distribution of layer-wise firing rates before/after alignment (from spike\_stats).}
  \caption{Per-layer mean firing rate across held-out subjects, with and without alignment.}
  \label{fig:rates}
\end{figure}
\begin{table}[t]
  \centering
  \caption{Single-factor ablations of S-EEGNeX-S (SHU 2022 pilot, split 41, learning rate $10^{-3}$; paired $\Delta$ accuracy vs the $10^{-3}$ S-EEGNeX-S, Wilcoxon). Neuron (H3) and encoding (H2) rows come from the main run at $10^{-4}$. \todo{Fill from outputs/analysis/snn\_round/shu\_pilot\_lr1e3/paired\_tests.csv (afe/spk/noinit) and the main SHU analysis (H2/H3).}}
  \label{tab:ablation}
  \small
  \begin{tabular}{lcc}
    \toprule
    Variant & $\Delta$ acc (pp) & $p$ \\
    \midrule
    Analog front-end (spikes after spatial filter) & \todo{$-1.82$} & \todo{0.027} \\
    Binary readout instead of LI & \todo{$+0.50$} & \todo{n.s.} \\
    Default BN init & \todo{$-0.30$} & \todo{n.s.} \\
    LIF / PLIF instead of HetPLIF (H3) & \todo{} & \todo{} \\
    Direct encoding, $T{=}4$ (H2) & \todo{} & \todo{} \\
    \bottomrule
  \end{tabular}
\end{table}
\todo{Energy: S-EEGNeX-S vs EEGNeX vs EEGNet (state EEGNeX accuracy at roughly EEGNet-scale energy; numbers from energy.csv; no multiplier in title/abstract).}

% --------------------------------------------------------------------------
\section{Conclusion}
\label{sec:conclusion}
\todo{One paragraph. Limitations: transductive alignment (unlabeled test-subject data), single training seed with shared initialisation, recipe dependence, SHU floor effect (chance band 51.4--53.7\%), temporal-native model $\sim$3$\times$ slower to train (launch-bound sequential loop).}

% Balance the last page if needed: \vfill\pagebreak

% --------------------------------------------------------------------------
% 5th page: references only (plus funding / compliance statement if required).
\bibliographystyle{IEEEbib}
\bibliography{refs}

\end{document}
